% Source: https://artofproblemsolving.com/community/c3h1533471_challenge_quadratics
% Stable identity: aops-topic-1533471-q01
After $k$ decreases of $\$10$ the price is $200-10k$ dollars and $90+5k$ jackets are sold, so the revenue is
\[
(200-10k)(90+5k)=18000+100k-50k^2 .
\]
\textbf{(a)} $18000+100k-50k^2=17600\iff k^2-2k-8=0\iff(k-4)(k+2)=0$. With $k\ge0$ we get $k=4$: the price is $\boxed{\$160}$ and $\boxed{110}$ jackets are sold ($160\cdot110=17600$). (The root $k=-2$ would mean raising the price to $\$220$ and selling $80$ jackets, which also gives $17600$, if the same rule applied to increases.)

\textbf{(b)} $18000+100k-50k^2\ge15600\iff k^2-2k-48\le0\iff(k-8)(k+6)\le0\iff-6\le k\le8$. The price $200-10k$ is lowest for the largest $k$, namely $k=8$: the lowest price is $\boxed{\$120}$, with $\boxed{130}$ jackets sold ($120\cdot130=15600$).
